\chapter{Filters, Matching Networks, and Transformers}
\label{ch:filters}
\section{Why Filters Matter in Amateur Radio}
Filters select desired signals and reject unwanted ones. They appear in receiver
front ends, intermediate-frequency (IF) stages, audio stages, transmitter
outputs, anti-aliasing paths,
power supplies, and interference cures. A filter is nothing more than a circuit
whose transfer function has been shaped on purpose---poles placed to form the
passband, zeros placed to notch the stopband---so everything from Part~II applies
directly.

\section{Basic Filter Types}
Four responses cover almost every filtering job, distinguished by which band of
frequencies they pass:

\begin{center}
\begin{tabularx}{\textwidth}{@{}L{0.26\textwidth}Y@{}}
\toprule
Type & What it does \\
\midrule
Low-pass & Passes low frequencies; attenuates high frequencies. \\
High-pass & Passes high frequencies; attenuates low frequencies. \\
Band-pass & Passes a band of frequencies; attenuates outside it. \\
Band-stop (notch) & Rejects a band of frequencies; passes outside it. \\
\bottomrule
\end{tabularx}
\end{center}

\section{Where Filter Behavior Comes From}
Two circuits already developed from first principles are the whole foundation of
this chapter, so we recall the results rather than re-deriving them:
\begin{description}
\item[First order --- the RC low-pass (\cref{ch:rc}).] One energy store, one pole
at \(s=-1/RC\), giving \(G(s)=1/(1+RCs)\): flat to the corner
\(\omega_c=1/RC\), then \(-\SI{20}{\decibel}\) per decade. Taken across the
resistor instead, the same circuit is a high-pass.
\item[Second order --- the LC low-pass (\cref{sec:lc-lowpass} in
\cref{ch:rlc}).] Two energy stores, a pole pair, giving
\[
G(s)=\frac{\omega_0^2}{s^2+\dfrac{\omega_0}{Q}s+\omega_0^2},
\qquad
\omega_0=\frac{1}{\sqrt{LC}},
\qquad
Q=\frac{1}{R}\sqrt{\frac{L}{C}},
\]
which falls at \(-\SI{40}{\decibel}\) per decade and peaks by \(\approx Q\) near
\(\omega_0\) when \(Q>1/\sqrt2\).
\end{description}
Those two supply everything a filter designer needs: a \emph{corner frequency}
set by the reactance values, a \emph{roll-off rate} set by the number of poles,
and---for second order and above---a \emph{peaking} set by the damping. The rest
of this chapter is about choosing them deliberately.

\section{Order, Energy Stores, and Roll-Off}
The pattern generalizes. Each independent energy-storage element (an \(L\) or a
\(C\)) contributes one pole, so an \(n\)th-order filter has \(n\) reactive
elements and rolls off at about \(-20n\,\si{\decibel}\) per decade far above
cutoff (each pole adding \(-\SI{20}{\decibel}\)/decade, each zero
\(+\SI{20}{\decibel}\)/decade). Steeper skirts require more poles, hence more
inductors and capacitors. The \emph{order} is thus both a circuit fact (how many
reactive parts) and a response fact (how fast it rolls off).

\section{Zeros in Radio Circuits: Traps and Notches}
If poles build the passband, \emph{zeros} are how a circuit deliberately
\emph{rejects} energy. A transmission zero on the \(\jj\omega\) axis forces the
output to zero at one frequency. Two dual constructions put one there. A series LC
branch to ground has zero impedance at its resonance, shorting that frequency to
ground and placing a transmission zero at \(\wzero=1/\sqrt{LC}\)---a trap. A
parallel LC branch in series with the signal path is the dual: it becomes an open
circuit at resonance, again blocking one frequency. Wherever the exam describes a
trap, a wave trap, or a notch rejecting a specific frequency, a zero has been
parked at that frequency (\cref{fig:pz-notch} in \cref{ch:splane} shows the
pole--zero picture). The elliptic family below uses exactly this---stopband
\(\jj\omega\)-axis zeros---to sharpen its skirts.

\section{Designing and Tuning a Filter}
Turn the derivation around and you have the design procedure. A filter is
specified by a target frequency response; designing it means \emph{placing the
poles and zeros} to meet that target and then realizing them with a circuit.
There are three independent choices:
\begin{description}
\item[Cutoff frequency] scale all \(L\) and \(C\) together: \(\wzero=1/\sqrt{LC}\)
shifts the whole response along the frequency axis without changing its shape.
\item[Order] the number of poles (reactive elements) sets skirt steepness,
\(-20n\,\si{\decibel}\) per decade.
\item[Shape (the ``family'')] for a given order, \emph{where} the poles sit trades
passband flatness against transition sharpness against phase linearity.
\end{description}

That last choice is what the named families are: different rules for placing the
poles (and zeros) of a given order. \Cref{fig:filter-families} shows them at the
same order, and \cref{fig:butterworth-poles} shows the archetype---Butterworth
poles equally spaced on a circle.

\begin{tabularx}{\textwidth}{@{}L{0.19\textwidth}Y@{}}
\toprule
Family & Pole/zero placement and consequence \\
\midrule
Butterworth & Poles equally spaced on a semicircle of radius \(\wzero\) in the left half-plane. Maximally flat passband, no ripple, moderate transition. \\
Chebyshev & Poles on an ellipse pushed closer to the \(\jj\omega\) axis. Sharper cutoff at the price of passband ripple and worse phase/transient behavior. \\
Elliptic (Cauer; ``elliptical'' in the pool) & Chebyshev-like poles plus transmission \emph{zeros} on the \(\jj\omega\) axis in the stopband. Sharpest transition; ripple in both bands. \\
Bessel & Poles positioned for maximally flat \emph{group delay} (linear phase); gentle amplitude cutoff. Good for pulse fidelity. \\
Crystal / cavity & Very high-\(Q\) resonators (poles hugging the axis) for narrow, highly selective band-pass shapes. \\
DSP FIR & Digital (DSP, digital signal processing) filter with no feedback (FIR, finite impulse response): no poles away from the origin, so it can be exactly linear phase. \\
DSP IIR & Recursive digital filter (IIR, infinite impulse response): has poles, so it mimics analog families efficiently but with nonlinear phase. \\
\bottomrule
\end{tabularx}

\begin{figure}[htbp]
\centering
\begin{tikzpicture}[>=Latex,scale=1.15]
\draw[->,black!70] (-2.7,0) -- (0.8,0) node[right]{\(\sigma\)};
\draw[->,black!70] (0,-2.0) -- (0,2.0) node[above]{\(\jj\omega\)};
\draw[GuideLine,dashed] (0,0) circle (1.6);
% 4th-order Butterworth poles: angles 112.5,157.5,202.5,247.5 deg
\foreach \a in {112.5,157.5,202.5,247.5}{
  \filldraw[GuideBlue] ({1.6*cos(\a)},{1.6*sin(\a)}) circle (2.2pt);
}
\node[GuideLine] at (1.15,1.35) {\small \(|s|=\wzero\)};
\node[GuideBlue] at (-2.15,1.2) {\small Butterworth};
\end{tikzpicture}
\caption{Fourth-order Butterworth poles: four left-half-plane poles equally
spaced on the circle of radius \(\wzero\). Higher order adds more poles on the
same circle; Chebyshev pulls them toward the \(\jj\omega\) axis for a sharper but
rippled response.}
\label{fig:butterworth-poles}
\end{figure}

The \emph{transition band} is the frequency span between the passband edge and the
stopband; a ``moderate'' transition (Butterworth) is wider than the ``sharp''
transition of a Chebyshev or elliptic filter of the same order.
\Cref{fig:filter-families} makes the flatness-versus-sharpness-versus-ripple trade
concrete.

\begin{figure}[htbp]
\centering
\includegraphics[width=0.9\textwidth]{filter_families}
\caption{Order-5 low-pass responses, all normalized to the same cutoff
\(\omega_c\). Butterworth is maximally flat with a moderate transition; Chebyshev
trades \SI{1}{\decibel} of passband ripple for a sharper skirt; elliptic is
sharpest of all, with ripple in both bands and stopband nulls (its
\(\jj\omega\)-axis zeros); Bessel is gentlest, chosen for clean phase and pulse
fidelity rather than sharpness.}
\label{fig:filter-families}
\end{figure}

\begin{advancedbox}
\textbf{Where those curves come from.} Each family is defined by a formula for the
magnitude response, from which the pole positions follow. For order \(n\) and cutoff
\(\omega_c\):
\[
|G|^2_{\text{Butterworth}}=\frac{1}{1+(\omega/\omega_c)^{2n}},
\qquad
|G|^2_{\text{Chebyshev}}=\frac{1}{1+\varepsilon^2T_n^2(\omega/\omega_c)},
\]
where \(T_n\) is the order-\(n\) Chebyshev polynomial and \(\varepsilon\) sets the
ripple. Butterworth's poles are the left-half-plane roots of
\(1+(s/\jj\omega_c)^{2n}=0\)---which is exactly why they land equally spaced on a
circle of radius \(\omega_c\) (\cref{fig:butterworth-poles}). Chebyshev's roots lie
on an ellipse; elliptic (Cauer) filters replace \(T_n\) with a rational elliptic
function, which is what introduces the finite stopband zeros; Bessel poles come from
Bessel polynomials chosen to flatten group delay. \Cref{fig:filter-families} was
generated numerically from these standard designs (via \texttt{scipy.signal}) rather
than sketched, so the ripple depths and skirt slopes are quantitatively right.
\end{advancedbox}

Finally, the same target poles can be \emph{realized} several ways: a passive LC
ladder, an active Sallen--Key stage (\cref{ch:active}) that makes complex poles
without inductors, a crystal or cavity when very high \(Q\) is needed, or a DSP
filter (FIR/IIR) that places the poles and zeros in software. Same poles,
different hardware---so ``customizing'' a filter means choosing the order (skirt
steepness), the family (the flatness/sharpness/phase trade), the cutoff (by
scaling \(L\) and \(C\)), and the realization that suits the frequency range and
available parts.

\section{Impedance Matching}
An impedance-matching network does two jobs:
\begin{enumerate}
\item cancel the reactive part of a complex load, and
\item transform the remaining resistance to the desired system impedance.
\end{enumerate}
In a \(\SI{50}{\ohm}\) RF system, matching maximizes power transfer and keeps
transmission-line reflections controlled.

\subsection{The L-Network}
The simplest matching network uses two reactances. To match two real resistances
\(R_{\text{hi}}>R_{\text{lo}}\), place a series reactance on the low-resistance
side and a shunt reactance on the high-resistance side. The network operating
\(Q\) is set by the resistance ratio:
\[
Q=\sqrt{\frac{R_{\text{hi}}}{R_{\text{lo}}}-1},
\qquad
X_{\text{series}}=Q\,R_{\text{lo}},
\qquad
X_{\text{shunt}}=\frac{R_{\text{hi}}}{Q}.
\]
The two reactances are of opposite type (one inductor, one capacitor); choosing
which is the inductor sets whether the network is low-pass or high-pass
(\cref{fig:lnetwork}). An L-network has no free parameter left over, so its
bandwidth is fixed once the resistances are chosen.

\begin{figure}[htbp]
\centering
\begin{circuitikz}[scale=1]
% source-side termination R_lo (drawn as a resistor to ground)
\draw (0,2) coordinate(lt) to[R,l_=\(R_{\text{lo}}\)] (0,0) coordinate(lb);
% series reactance on the low-resistance side
\draw (lt) to[L,l=\(X_{\text{series}}\)] (3,2) coordinate(mid);
% shunt reactance on the high-resistance side
\draw (mid) to[C,l=\(X_{\text{shunt}}\)] (3,0) coordinate(mb);
% load-side termination R_hi (drawn as a resistor to ground)
\draw (mid) to[short] (5.5,2) coordinate(rt);
\draw (rt) to[R,l=\(R_{\text{hi}}\)] (5.5,0) coordinate(rb);
% ground rail
\draw (lb) -- (rb);
\node[ground] at (3,0){};
% input / output port labels
\node[above=0.05cm,GuideBlue] at (lt) {\small input (source side)};
\node[above=0.05cm,GuideBlue] at (rt) {\small output (load side)};
\draw[GuideBlue,-{Latex}] (-0.55,1.3) -- (-0.55,1.85);
\draw[GuideBlue,-{Latex}] (6.05,1.85) -- (6.05,1.3);
\end{circuitikz}
\caption{An L-network matching a low resistance \(R_{\text{lo}}\) to a high one
\(R_{\text{hi}}\) (both drawn as resistors to the ground rail). The network sits
\emph{between} the two terminations: the source drives the left-hand port and the
load hangs on the right-hand one. The series reactance goes on the low-resistance
side, the shunt reactance on the high-resistance side. Swapping which element is the
inductor makes it high-pass instead of low-pass.}
\label{fig:lnetwork}
\end{figure}

\subsection{Pi and T Networks}
Adding a third element gives a \(\pi\) or T network (\cref{fig:pit-network}),
which can be viewed as two L-networks back to back that meet at an internal
\emph{virtual resistance} \(R_v\).

\begin{figure}[htbp]
\centering
\begin{circuitikz}[scale=0.9]
% --- Pi network ---
\draw (0,2) coordinate(pin) to[L,l=\(L\)] (3,2) coordinate(pout);
\draw (pin) to[C,l_=\(C_1\)] (0,0) coordinate(pinb);
\draw (pout) to[C,l=\(C_2\)] (3,0) coordinate(poutb);
\draw (pinb) -- (poutb);
\node[ground] at (1.5,0){};
\draw (pin) to[short,-o] (-0.9,2) node[left]{in};
\draw (pout) to[short,-o] (3.9,2) node[right]{out};
\node at (1.5,-1.0) {\(\pi\) network};
\end{circuitikz}
\qquad
\begin{circuitikz}[scale=0.9]
% --- T network ---
\draw (0,2) coordinate(tin) to[L,l=\(L_1\)] (2,2) coordinate(tm)
      to[L,l=\(L_2\)] (4,2) coordinate(tout);
\draw (tm) to[C,l=\(C\)] (2,0) coordinate(tmb);
\draw (0,0) coordinate(tinb) -- (4,0) coordinate(toutb);
\node[ground] at (2,0){};
\draw (tin) to[short,-o] (-0.9,2) node[left]{in};
\draw (tout) to[short,-o] (4.9,2) node[right]{out};
\node at (2,-1.0) {T network};
\end{circuitikz}
\caption{Three-element matching networks. The \(\pi\) network (two shunt
capacitors, one series inductor) and the T network (two series inductors, one
shunt capacitor) each add a third element over the L-network, giving the extra
degree of freedom that sets the loaded \(Q\) through the internal virtual
resistance \(R_v\).}
\label{fig:pit-network}
\end{figure}

Choosing \(R_v\) below both terminations sets a loaded \(Q\) of roughly
\[
Q_{\text{loaded}}\approx\sqrt{\frac{R_{\max}}{R_v}-1},
\]
so \(R_v\) is the extra degree of freedom that lets the designer pick the loaded
\(Q\)---and therefore the bandwidth and harmonic rejection---independently of the
resistance ratio. This is why transmitter output and antenna-tuner networks are
usually \(\pi\) or T: they match \emph{and} set selectivity. Higher loaded \(Q\)
narrows bandwidth and improves harmonic suppression at the cost of higher
circulating currents and losses.

\begin{mistakebox}
Matching is not the same thing as filtering. A network can do both, but the
concepts are distinct: matching controls the impedance seen by a source or line,
while filtering controls frequency-dependent transmission. A \(\pi\)-network tuner
happens to be a low-pass filter too, which is a feature, not the definition.
\end{mistakebox}

\section{Transformers}
For an ideal transformer (\cref{fig:transformer}),
\[
\frac{V_s}{V_p}=\frac{N_s}{N_p},
\qquad
\frac{I_s}{I_p}=\frac{N_p}{N_s},
\qquad
Z_p=\left(\frac{N_p}{N_s}\right)^2 Z_s.
\]

\begin{figure}[htbp]
\centering
\begin{circuitikz}[scale=1]
\node[transformer core] (T) at (0,0){};
\draw (T.A1) to[short,-o] ++(-1.1,0) node[left]{\(+\)} node[left=0.7cm]{\(V_p\)};
\draw (T.A2) to[short,-o] ++(-1.1,0) node[left]{\(-\)};
\draw (T.B1) to[short,-o] ++(1.1,0) node[right]{\(+\)} node[right=0.7cm]{\(V_s\)};
\draw (T.B2) to[short,-o] ++(1.1,0) node[right]{\(-\)};
\node[above=0.15cm] at (T.A1) {\(N_p\)};
\node[above=0.15cm] at (T.B1) {\(N_s\)};
\end{circuitikz}
\caption{Ideal two-winding transformer. The primary (\(N_p\) turns) and secondary
(\(N_s\) turns) share a common core; voltages scale as the turns ratio
\(N_s/N_p\), currents inversely, and impedance as the ratio squared.}
\label{fig:transformer}
\end{figure}
Impedance transforms as the turns ratio \emph{squared}. \textbf{Core saturation}
should be avoided: every magnetic core has a maximum flux density it can carry,
and once the drive pushes it near that limit the core's permeability collapses, so
the inductance drops and the winding current shoots up. The result is distortion,
heating, and loss of the intended impedance transformation---the reason
transformers and inductors are rated for a maximum voltage-time product (or power)
at a given frequency.

\section{Worked Example: Quarter-Wave and L-Network}
\textbf{Problem.} A transmitter feeds a \(\SI{50}{\ohm}\) system, but the load
(say an antenna at resonance) presents a purely resistive \(\SI{200}{\ohm}\).
Design two matching networks that make the \(\SI{200}{\ohm}\) load look like
\(\SI{50}{\ohm}\) to the source: (a) a quarter-wave transformer, and (b) an
L-network. Compare them.

To match a real \(\SI{200}{\ohm}\) load to a \(\SI{50}{\ohm}\) system with a
lossless quarter-wave transformer (the formula \(Z_t=\sqrt{Z_0Z_L}\) is derived in
\cref{ch:lines}),
\[
Z_t=\sqrt{Z_0Z_L}=\sqrt{(50)(200)}=\SI{100}{\ohm}.
\]
The same match with an L-network gives
\[
Q=\sqrt{\frac{200}{50}-1}=\sqrt{3}\approx1.73,
\qquad
X_{\text{series}}=1.73\times50\approx\SI{87}{\ohm},
\qquad
X_{\text{shunt}}=\frac{200}{1.73}\approx\SI{116}{\ohm}.
\]
Both realize the match; the quarter-wave version needs a line one quarter
wavelength long at the design frequency, while the L-network is lumped and
narrower band.
